% =============================================================================
% B.Comp. Dissertation (CP4101) -- Machine Learning in Prediction Markets
% Format follows "CP4101 B.Comp. Dissertation Format" (May 2008):
%   - A4, Times 12pt, 1.5 line spacing for the main report
%   - Main report <= 55 pages (Introduction -> Conclusion), numbered 1,2,3,...
%   - Front matter numbered i, ii, ...; references continue main numbering
%   - Order: cover, title page, abstract, acknowledgement, contents, lists,
%            main report, references, appendices
%   - Citations: author-year (APA style); references sorted by first author
%   - Optional: <=10 page conference paper in the appendix (on top of 55)
% Everything marked \todo{...} is a placeholder to be replaced.
% Build: pdflatex main && pdflatex main   (no BibTeX needed; refs are inline)
% =============================================================================
\documentclass[12pt,a4paper,oneside]{report}

\usepackage[T1]{fontenc}
\usepackage{mathptmx}                 % Times Roman text and math
\usepackage[margin=2.54cm]{geometry}
\usepackage{setspace}
\usepackage{amsmath,amssymb}
\usepackage{graphicx}
\usepackage{booktabs}
\usepackage{xcolor}
\usepackage{enumitem}
\usepackage{caption}
\usepackage[authoryear,round]{natbib}
\usepackage[hidelinks]{hyperref}
\usepackage{titlesec}
\usepackage{fancyhdr}

\onehalfspacing
\setlength{\parskip}{4pt}
\setlist{nosep}
\captionsetup{font=small}

% Keep heading depth shallow (format asks for minimal sub-heading depth)
\setcounter{secnumdepth}{2}
\setcounter{tocdepth}{2}

% ---- Draft helpers ----------------------------------------------------------
\newcommand{\todo}[1]{\textcolor{red}{[TODO: #1]}}
\newcommand{\figplaceholder}[2]{%
  \begin{figure}[t]\centering
    \fbox{\parbox{0.8\textwidth}{\centering\vspace{2em}\todo{#1}\vspace{2em}}}
    \caption{#2}
  \end{figure}}

% ---- Project metadata (edit here) -------------------------------------------
\newcommand{\ProjTitle}{Machine Learning in Prediction Markets}
\newcommand{\ProjAuthor}{Casper Ng}          % \todo: full name as on matriculation
\newcommand{\ProjYear}{2026/2027}
\newcommand{\ProjNo}{\todo{Project No, e.g. H123456}}
\newcommand{\ProjAdvisor}{\todo{rank, e.g. Assoc Prof} Lu Yao}

\begin{document}

% =============================================================================
% FRONT COVER (hard cover)
% =============================================================================
\pagestyle{empty}
\begin{titlepage}
\centering
\vspace*{3cm}
{\large B.Comp. Dissertation}\\[3.5cm]
{\LARGE\bfseries \ProjTitle}\\[3.5cm]
{\large By}\\[0.6cm]
{\large \ProjAuthor}\\[4cm]
{\large Department of Computer Science}\\[0.3cm]
{\large School of Computing}\\[0.3cm]
{\large National University of Singapore}\\[0.6cm]
{\large \ProjYear}
\end{titlepage}

% =============================================================================
% TITLE PAGE (adds project number, advisor, deliverables)
% =============================================================================
\begin{titlepage}
\centering
\vspace*{1.5cm}
{\large B.Comp. Dissertation}\\[1.5cm]
{\LARGE\bfseries \ProjTitle}\\[2cm]
{\large By}\\[0.5cm]
{\large \ProjAuthor}\\[2cm]
{\large Department of Computer Science}\\[0.3cm]
{\large School of Computing}\\[0.3cm]
{\large National University of Singapore}\\[0.5cm]
{\large \ProjYear}\\[2cm]
\begin{flushleft}
\hspace{1.5cm}Project No: \ProjNo\\
\hspace{1.5cm}Advisor: \ProjAdvisor\\[0.6cm]
\hspace{1.5cm}Deliverables:\\
\hspace{3cm}Report: 1 Volume\\
\hspace{3cm}Source code and data pipeline: \todo{repository / archive}\\
\end{flushleft}
\end{titlepage}

% =============================================================================
% FRONT MATTER (roman numerals)
% =============================================================================
\pagestyle{plain}
\pagenumbering{roman}
\setcounter{page}{1}

% ---- Abstract (<= 200 words; <= 5 subject descriptors; <= 5 keywords) --------
\chapter*{Abstract}
\addcontentsline{toc}{chapter}{Abstract}
\singlespacing
\todo{Keep to $\le$ 200 words. Draft below; rewrite once results are final.}

Prediction markets such as Kalshi let people trade contracts that pay out if a
real-world event happens, so their prices can be read as probabilities. This
project asks whether machine learning can forecast short-term price movements
in these markets better than the market itself, after realistic trading costs.
We study two settings. First, we test whether the outcome of one scheduled
economic release (for example inflation or jobs data) predicts the repricing of
other related contracts, using a spatial-temporal graph model with a learned
graph structure. Second, we build a forecaster for live sports markets that
turns play-by-play data into a price prediction for the moment a play happens,
and evaluate it with a strict as-of protocol, a frozen holdout, and a live
paper-trading deployment. \todo{One sentence on headline result, e.g. what
worked, what did not, and the net-of-cost conclusion.}

\vspace{0.6em}
\noindent\textbf{Subject Descriptors:}\\[2pt]
\hspace*{1.2cm}I.2.6 Learning\\
\hspace*{1.2cm}I.5.1 Models\\
\hspace*{1.2cm}G.3 Probability and Statistics\\
\hspace*{1.2cm}J.4 Social and Behavioral Sciences (Economics)\\
\hspace*{1.2cm}\todo{verify codes against the current ACM CCS / Computing Reviews scheme with advisor}

\vspace{0.6em}
\noindent\textbf{Keywords:}\\[2pt]
\hspace*{1.2cm}prediction markets, graph neural networks, time series forecasting,
market efficiency, sequential decision making

\vspace{0.6em}
\noindent\textbf{Implementation Software and Hardware:}\\[2pt]
\hspace*{1.2cm}Python 3, PyTorch, pandas, Slurm GPU cluster (NVIDIA A100),
\todo{OS, cloud sandbox, library versions}
\onehalfspacing

% ---- Acknowledgement (one double-spaced page) --------------------------------
\chapter*{Acknowledgement}
\addcontentsline{toc}{chapter}{Acknowledgement}
\doublespacing
\todo{One page. Thank advisor Lu Yao, the evaluator, the School of Computing for GPU cluster access, family and friends.}
\onehalfspacing

% ---- Contents and lists ------------------------------------------------------
\singlespacing
\tableofcontents
\clearpage
\addcontentsline{toc}{chapter}{List of Figures}
\listoffigures
\clearpage
\addcontentsline{toc}{chapter}{List of Tables}
\listoftables
\clearpage
\chapter*{List of Abbreviations}
\addcontentsline{toc}{chapter}{List of Abbreviations}
\begin{tabular}{@{}ll@{}}
AGCRN & Adaptive Graph Convolutional Recurrent Network\\
GRU   & Gated Recurrent Unit\\
STG   & Spatial-Temporal Graph\\
GNN   & Graph Neural Network\\
OOS   & Out of sample\\
PBP   & Play-by-play (event feed of a game)\\
\todo{add as needed} &
\end{tabular}
\onehalfspacing

% =============================================================================
% MAIN REPORT (arabic numerals, starts at 1; max 55 pages through Conclusion)
% =============================================================================
\clearpage
\pagenumbering{arabic}
\setcounter{page}{1}

% -----------------------------------------------------------------------------
\chapter{Introduction}
\label{ch:intro}
% -----------------------------------------------------------------------------

\section{Motivation}
Prediction markets are exchanges where each contract pays \$1 if a stated event
occurs and \$0 otherwise, so a price of 63 cents is read as a 63\% chance
\citep{wolfers2004prediction,arrow2008promise}. Regulated venues such as Kalshi
now list thousands of contracts on economic data, policy decisions and sports.
\todo{Why this matters for ML: high-frequency, labelled, binary-outcome data with an observable ground truth; a clean testbed for forecasting under distribution shift and transaction costs.}

\section{Problem Statement}
\todo{State the central question in plain words, e.g.: ``Can a learned model
predict short-horizon price changes in prediction markets well enough to beat
the market after trading costs, and how should such a claim be tested honestly?''}
Sub-questions:
\begin{enumerate}
  \item Does information released in one market (a macro data release) move the
        prices of related markets in a predictable way (Chapter~\ref{ch:macro})?
  \item Can a sequence/graph model of live game events forecast sports-contract
        prices ahead of the market (Chapter~\ref{ch:ingame})?
  \item Do such forecasts survive realistic execution, fees and a live
        forward test (Chapter~\ref{ch:trading})?
\end{enumerate}

\section{Contributions}
\todo{3--5 bullets. Candidate list based on the repo:}
\begin{itemize}
  \item A strategy-pattern framework for building spatial-temporal graph
        datasets from event-driven market data.
  \item A negative-result study showing when learned graph structure
        (AGCRN) does not beat imposed economic structure or linear baselines,
        including a semi-synthetic structure-recovery test.
  \item A strict as-of evaluation protocol that avoids look-ahead in features,
        baselines and instrument selection.
  \item An in-game forecasting model with a frozen holdout specification and
        a live paper-trading system.
  \item An analysis of execution cost, adverse selection and feed latency.
\end{itemize}

\section{Report Outline}
\todo{One short paragraph mapping chapters to the questions above.}

% -----------------------------------------------------------------------------
\chapter{Background and Related Work}
\label{ch:background}
% -----------------------------------------------------------------------------
\section{Prediction Markets and Kalshi}
Contract structure, order book, taker/maker, fee schedule (fee grows with
$p(1-p)$), settlement. \todo{Cite Kalshi documentation; keep to what the
models need.}

\section{Market Efficiency and Underreaction}
Efficient-markets framing; evidence for slow information absorption
\citep{hong2000bad}; microstructure and signed order flow \citep{lee1991inferring}.
\todo{Add prediction-market-specific efficiency / favourite--longshot bias
literature; add the sports-market paper discussed in the research summary.}

\section{Learning on Time Series and Graphs}
\begin{itemize}
  \item Recurrent models: GRU \citep{cho2014learning}.
  \item Spatial-temporal graph networks and adaptive graph learning:
        AGCRN \citep{bai2020adaptive}. \todo{Add Graph WaveNet, DCRNN, STGCN.}
  \item Gradient-boosted trees and linear baselines as strong references.
\end{itemize}

\section{Evaluating ML Forecasts in Finance}
Look-ahead bias, overlapping labels, multiple testing, and why published gains
often disappear out of sample \citep{bailey2014pseudo,lopezdeprado2018advances}.
\todo{Link to literature notes in reports/literature\_stg\_finance.md.}

\section{Machine Learning for Sports Betting}
Surveys cover the algorithms and evaluation practices used
\citep{bunker2019machine,galekwa2024systematic}. Two results shape our design:
choosing models by calibration rather than accuracy gives better betting returns
\citep{walsh2024machine}, and a model should be rewarded for differing from the
bookmaker's odds rather than for matching them \citep{hubacek2019exploiting}.
For in-game forecasting, simple calibrated models are a strong baseline
\citep{pelechrinis2017iwinrnfl}. \todo{Add the NBA recurrent-model thesis (Poole)
and discuss that prior work targets bookmaker odds, whereas we target
prediction-market prices with an explicit execution model.}

\section{Positioning of This Work}
\todo{Short comparison table: prior work vs.\ this project (data, task, evaluation, costs).}

% -----------------------------------------------------------------------------
\chapter{Data and Problem Formulation}
\label{ch:data}
% -----------------------------------------------------------------------------
\section{Data Sources}
\begin{itemize}
  \item Kalshi trade archive (tens of millions of trades; labelled taker side)
        and market metadata. \todo{final counts and date range}
  \item Macro release calendar with consensus forecasts, and macro indicators.
  \item Play-by-play feeds for MLB, NBA, NFL. \todo{sources, latency table}
\end{itemize}
\todo{Table: dataset, size, period, source, known gaps (e.g.\ no historical quotes, archive concentrated in 2025).}

\section{Strict As-Of Principle}
Every feature and every normaliser is computed only from information available
at the decision time; baselines are rolling or expanding, never full-history.
\todo{Describe the as-of layer and the integrity checks.}

\section{Task Definitions}
\todo{Formalise both tasks: inputs, target (price change over horizon $h$ or
settlement value), decision time, metrics (R$^2$, AUC, gross and net P\&L per
contract).}

\section{Splits and Evaluation Protocol}
Walk-forward / time-ordered splits, block permutation tests, frozen holdout.
\todo{Describe split dates and what is frozen before looking at results.}

% -----------------------------------------------------------------------------
\chapter{Study I: Cross-Market Propagation in Macro Markets}
\label{ch:macro}
% -----------------------------------------------------------------------------
\section{Spatial-Temporal Graph Framework}
Nodes, edges, snapshots and labels as interchangeable strategies; resolution-triggered
graph snapshots. \todo{Figure: pipeline diagram.}

\section{Models}
\subsection{Imposed-Structure Baselines}
Zero-parameter economic sign rule; per-pair regression; linear ladder.
\subsection{Learned Graph Models}
Inductive AGCRN with adaptive adjacency $\tilde A=\mathrm{softmax}(\mathrm{ReLU}(EE^\top))$
and node-adaptive weights \citep{bai2020adaptive}. \todo{Equations, hyper-parameters.}
\subsection{Other Learners}
Tree ensembles, MLP, Bayesian hierarchical model.

\section{Results}
\todo{Tables: model comparison on R$^2$/AUC; structure estimation (edges surviving
multiple-testing correction); gross edge versus cost.}

\section{Can a Graph Be Learned From This Data?}
Semi-synthetic recovery test: plant a known graph in data built from the real
release calendar, and measure whether each learner recovers it
(Section: \todo{numbers from reports/recovery\_test.md}).

\section{Lessons}
Why the learned graph did not help (too few co-firing events, signal dilution,
implementation pitfalls found by checklist), and why this motivated moving to a
data-rich setting.

% -----------------------------------------------------------------------------
\chapter{Study II: In-Game Sports Forecasting}
\label{ch:ingame}
% -----------------------------------------------------------------------------
\section{Problem Setting}
Given the live feed of a game, predict the price of related Kalshi contracts a
short time after each play. \todo{Plain-language definition of a play, a
contract leg, and the decision time.}

\section{Data Pipeline}
Play-by-play alignment to market time, feed latency and revisions, as-of
feature construction. \todo{Figure: timeline of event $\to$ feed $\to$ price.}

\section{Model Architecture}
GRU sequence encoder over plays with a graph over related contracts within a game.
\todo{Architecture figure, input features, loss, training on GPU cluster, seeds
(5-seed average).}

\section{Ablations}
Feed delay, graph vs.\ no graph, feature groups, per-sport differences.
\todo{Ablation table.}

\section{Frozen Holdout Results}
Report what the frozen specification (candidates, exit rules, sizing) produced,
with no post-hoc pass rule. \todo{Gross P\&L per contract with confidence intervals.}

% -----------------------------------------------------------------------------
\chapter{Trading System and Live Evaluation}
\label{ch:trading}
% -----------------------------------------------------------------------------
\section{Execution Model}
Maker vs.\ taker, fee model, fill assumptions, order-management variants.
\todo{Gross first, then fees.}

\section{Why Signals Shrink: Stale Prices and Adverse Selection}
Stale last print versus true move; settled losers as post-event pick-offs.
\todo{Tables from reports/ingame\_journal.md and stale-reference analysis.}

\section{Live Paper-Trading System}
Runners, recorder, replay tool, watchdog, sandbox deployment.
\todo{System diagram.} \todo{Live vs.\ backtest parity results from forward\_test\_log.md.}

\section{Forward-Test Results}
\todo{Games traded, P\&L, divergences from backtest (feed latency, postponements).}

% -----------------------------------------------------------------------------
\chapter{Discussion}
\label{ch:discussion}
% -----------------------------------------------------------------------------
\section{What Worked and What Did Not}
\todo{Honest critique: where ML helped, where simple structure won, data limits.}

\section{Threats to Validity}
Look-ahead leakage risks, multiple testing, regime change, archive selection,
feed latency in live deployment.

\section{Limitations}
\todo{Scope: MLB/NBA/NFL only, paper trading, capacity and market impact.}

% -----------------------------------------------------------------------------
\chapter{Conclusion and Future Work}
\label{ch:conclusion}
% -----------------------------------------------------------------------------
\section{Summary}
\todo{Answer each sub-question from Section 1.2 in one paragraph.}

\section{Recommendations for Further Work}
\todo{Deep-book/quote data, lower-latency feeds, additional sports, calibration
of probabilities, online learning.}

% =============================================================================
% REFERENCES (page numbering continues; alphabetical by first author)
% =============================================================================
\singlespacing
\renewcommand{\bibname}{References}
\begin{thebibliography}{99}
\addcontentsline{toc}{chapter}{References}
\setlength{\itemsep}{4pt}

\bibitem[Arrow et al.(2008)]{arrow2008promise}
Arrow, K.J., Forsythe, R., Gorham, M., Hahn, R., Hanson, R., Ledyard, J.O.,
Levmore, S., Litan, R., Milgrom, P., Nelson, F.D., Neumann, G.R., Ottaviani, M.,
Schelling, T.C., Shiller, R.J., Smith, V.L., Snowberg, E., Sunstein, C.R.,
Tetlock, P.C., Tetlock, P.E., Varian, H.R., Wolfers, J. and Zitzewitz, E. (2008).
The promise of prediction markets. \textit{Science}, Vol.~320, No.~5878,
May 2008, pp.~877--878.

\bibitem[Bai et al.(2020)]{bai2020adaptive}
Bai, L., Yao, L., Li, C., Wang, X. and Wang, C. (2020). Adaptive graph
convolutional recurrent network for traffic forecasting. In
\textit{Advances in Neural Information Processing Systems 33 (NeurIPS 2020)}.

\bibitem[Bailey et al.(2014)]{bailey2014pseudo}
Bailey, D.H., Borwein, J.M., L\'opez de Prado, M. and Zhu, Q.J. (2014).
Pseudo-mathematics and financial charlatanism: the effects of backtest
overfitting on out-of-sample performance. \textit{Notices of the American
Mathematical Society}, Vol.~61, No.~5, May 2014, pp.~458--471.

\bibitem[Bunker and Susnjak(2019)]{bunker2019machine}
Bunker, R.P. and Susnjak, T. (2019). The application of machine learning
techniques for predicting results in team sport: a review. arXiv:1912.11762.
\todo{check final published venue and year}

\bibitem[Cho et al.(2014)]{cho2014learning}
Cho, K., van Merri\"enboer, B., G\"ul\c{c}ehre, \c{C}., Bahdanau, D.,
Bougares, F., Schwenk, H. and Bengio, Y. (2014). Learning phrase representations
using RNN encoder--decoder for statistical machine translation. In
\textit{Proceedings of EMNLP 2014}, pp.~1724--1734.

\bibitem[Galekwa et al.(2024)]{galekwa2024systematic}
Galekwa, R.M., Tshimula, J.M., Tajeuna, E.G. and Kyamakya, K. (2024).
A systematic review of machine learning in sports betting: techniques,
challenges, and future directions. arXiv:2410.21484.
\todo{verify author list and venue}

\bibitem[Hong et al.(2000)]{hong2000bad}
Hong, H., Lim, T. and Stein, J.C. (2000). Bad news travels slowly: size,
analyst coverage, and the profitability of momentum strategies.
\textit{The Journal of Finance}, Vol.~55, No.~1, February 2000, pp.~265--295.

\bibitem[Hubáček et al.(2019)]{hubacek2019exploiting}
Hub\'a\v{c}ek, O., \v{S}ourek, G. and \v{Z}elezn\'y, F. (2019). Exploiting
sports-betting market using machine learning. \textit{International Journal
of Forecasting}, Vol.~35, No.~2, 2019, pp.~783--796. \todo{verify volume/pages}

\bibitem[Lee and Ready(1991)]{lee1991inferring}
Lee, C.M.C. and Ready, M.J. (1991). Inferring trade direction from intraday
data. \textit{The Journal of Finance}, Vol.~46, No.~2, June 1991, pp.~733--746.

\bibitem[L\'opez de Prado(2018)]{lopezdeprado2018advances}
L\'opez de Prado, M. (2018). \textit{Advances in Financial Machine Learning}.
Wiley, Hoboken, NJ, 2018.

\bibitem[Pelechrinis(2017)]{pelechrinis2017iwinrnfl}
Pelechrinis, K. (2017). iWinRNFL: a simple, interpretable and well-calibrated
in-game win probability model for NFL. arXiv:1704.00197.

\bibitem[Walsh and Joshi(2024)]{walsh2024machine}
Walsh, C. and Joshi, A. (2024). Machine learning for sports betting: should
model selection be based on accuracy or calibration? \textit{Machine Learning
with Applications}, Vol.~16, 2024. \todo{verify article number}

\bibitem[Wolfers and Zitzewitz(2004)]{wolfers2004prediction}
Wolfers, J. and Zitzewitz, E. (2004). Prediction markets. \textit{Journal of
Economic Perspectives}, Vol.~18, No.~2, Spring 2004, pp.~107--126.

% \todo{Add: Kalshi documentation (organisation as author), Graph WaveNet,
% DCRNN, STGCN, the sports-market efficiency paper, the five papers in papers/.
% Verify every entry above before submission.}
\end{thebibliography}

% =============================================================================
% APPENDICES (single spacing, may use smaller font)
% =============================================================================
\appendix
\small
\singlespacing

\chapter{Implementation Details}
\todo{Repository layout, key modules, hyper-parameters, cluster job settings.}

\chapter{Additional Results}
\todo{Extra tables and plots that would break the flow of the main report.}

\chapter{Conference Paper (optional, $\le$10 pages)}
\todo{Delete this appendix if not submitting a paper. Allowed on top of the 55-page limit.}

\chapter{Reproducibility and Data Dictionary}
\todo{Data schemas, as-of layer definitions, how to rerun experiments.}

\end{document}
